\section{Introduction}\label{sec:intro}

\paragraph{A puzzle from practice.}
Data assimilation (DA) combines a dynamical model with observations to estimate
the state of a system~\cite{carrassi2018da}. Its premise is that the model
carries information the observations lack, so that a model-based estimate should
beat any estimate that ignores the model. Practice keeps contradicting this
premise. Gridded sea-surface height from observation-only optimal
interpolation resolves finer scales than an eddy-permitting ocean reanalysis
over the Gulf Stream (an effective resolution of about 151 against 239\,km)
\todo{cite DUACS, GLORYS; effective-resolution reference}. In a
controlled setting, a variational scheme with a learned prior reached a lower
error than the same scheme run with the true equations of
motion~\cite{fablet4dvarnet}. Learned and hybrid schemes are now routinely
reported to beat classical DA on the same data
\cite{brajard2020combining,bocquet2020bayesian,rozet2023sda}.

Three explanations are usually offered: the dynamical model is wrong; the DA
algorithm loses information (filtering instead of smoothing, Gaussian updates,
finite ensembles, local optimisation); or the scheme does not use all the
information available to it. They are rarely separated, because each scheme
mixes all three, and a comparison of RMSEs cannot tell them apart.

\figplaceholder{Fig.~1 -- the puzzle}{(a) Sea-surface height over the Gulf Stream: observation-only product vs model-based reanalysis. (b) A learned prior vs the true model in a variational scheme \cite{fablet4dvarnet}.}

\paragraph{The idea.}
We treat every DA scheme, classical or learned, as an estimator of the
conditional distribution $p(\xv \mid C)$ of the state given its context, and of
its mean $\E[\xv \mid C]$. The error of a scheme then splits along a chain of
idealisations (Sect.~\ref{sec:chain}): an irreducible term; a
\emph{restriction}, because the scheme uses less information than $C$; a
\emph{misspecification}, because it believes a model other than the true one;
and an \emph{approximation}, because it computes its own target only
approximately. Two schemes with the same RMSE can have very different budgets,
and only the budget says what to fix.

The same chain applies to the uncertainty a scheme reports. We score it with a
flow-matching score, $\FMS_\tau$, which equals the mean-squared error at $\tau =
0$ and is strictly proper for $\tau \in (0,1)$ (Sect.~\ref{sec:fms}). The split
into an irreducible part, a restriction part and a remainder is exact at every
$\tau$. A restriction raises the score curve; a misspecification bends it upward
at large $\tau$, where calibration dominates. This gives the paper's thesis an
operational form: \emph{a restriction costs variance honestly and preserves
calibration; a misspecification costs bias silently and destroys it}.

\paragraph{Questions.}
We ask three questions on a two-scale Lorenz-96 system, with a perfect model and
with a biased one, and with five classical and five learned schemes under the
same observations and the same tuning protocol.
\begin{enumerate}[label=\textbf{Q\arabic*},leftmargin=*]
  \item With a perfect model, how close does each scheme get to $\E[\xv\mid C]$,
    and where does the rest of its error go? (Sect.~\ref{sec:q1})
  \item What does model error do to the budget, and to the uncertainty each
    scheme reports? (Sect.~\ref{sec:q2})
  \item Does the form of the prior --- derived from the model equations or
    learned from data --- drive performance? (Sect.~\ref{sec:q3})
\end{enumerate}

\paragraph{Contributions.}
\begin{itemize}[leftmargin=*]
  \item The error chain of the conditional expectation, applied to classical and
    learned DA under matched information, with each of its terms estimated by a
    controlled contrast: smoothing, ensemble size, optimisation budget,
    training budget, conditioning and model error.
  \item Its extension along $\tau$ with $\FMS_\tau$, a proper score computable
    for point estimates, Gaussians and samples alike, with an exact split at
    every $\tau$.
  \item A controlled study on a two-scale Lorenz-96 system: a perfect model and
    model error, a regular and a random observing system, and every inference
    hyper-parameter tuned on validation windows.
\end{itemize}

\paragraph{What this paper is not.}
It proposes no new DA method; the learned schemes are instruments that relax one
commitment at a time. The target is the state only, not the parameters. Unrolled
variational solvers are left to a companion study \todo{P2}, and so is the general
representation of DA algorithms as flows \todo{ML paper}. One idealised system
is the price of matched information; we discuss transfer in
Sect.~\ref{sec:discussion}.
